\documentclass[10pt,a4paper]{amsart}
\usepackage[margin=19mm]{geometry}
\usepackage{amsmath,amssymb,mathtools}
\usepackage[scr=rsfs,cal=euler]{mathalfa}
\usepackage{xfrac}
\usepackage[unicode,hidelinks,bookmarks=false]{hyperref}
\newcommand{\ie}{i.e.\ }
\newcommand{\cf}{cf.\ }
\newcommand{\resp}{respectively\ }
\newcommand{\Subsection}{\subsection}
\providecommand{\nobreakdash}{\nobreak\hbox{-}}
\setlength{\emergencystretch}{2em}
\multlinegap0pt
\usepackage{amssymb}
\newtheorem{theorem}{Theorem}
\title{A $q$-microscope for mock-theta denominator patterns}
\author{Mohamed El Bachraoui}
\date{Source: arXiv:2609.05961v1 (CC BY 4.0)}
\begin{document}
\maketitle

\noindent\emph{Source and licence.} A $q$-microscope for mock-theta denominator patterns. Version arXiv:2609.05961v1 (CC BY 4.0), distributed by arXiv under the Creative Commons Attribution 4.0 International licence, http://creativecommons.org/licenses/by/4.0/, as recorded in the arXiv licence metadata for this exact version and shown on the abstract page https://arxiv.org/abs/2609.05961v1. Full licence notice: \texttt{licenses/arxiv-2609.05961-CC-BY-4.0.txt}.\par\smallskip

\noindent\textit{Editorial scope.} Counted unit: the article's theorem thm:nu-finite-theta (source0.tex lines 170--183). The article's complete original proof of it is delivered with it (source0.tex lines 553--677, one \texttt{proof} environment of 125 lines, which the article places away from the statement). No mathematics is written, repaired or re-derived: the statement and the proof are the article's own lines, in the article's own notation, and no retained line is substituted or re-flowed.  No further source-local context is needed: every cross-reference of every retained line resolves inside the two delivered spans.  Nothing was omitted from any retained span, so the package carries no omission record.  No retained line contains a citation, so the wrapper carries no bibliography and no bibliography entry.  No retained line contains a figure environment, an image-inclusion command or drawing code, so no image is delivered and the package declares no asset dependency.  Editorial formatting only, in addition: an AMS \texttt{amsart} preamble that restates the article's own declarations and macros used by the retained lines, namely the environments \texttt{theorem} on separate counters, together with the standard packages those lines need.  The retained environments print in this document as Theorem 1 in order of appearance, whereas the article prints the same items with the numbering of its own full document.  The complete native source of the article as distributed in the arXiv e-print for this exact version is bundled unchanged as \texttt{sources/209479/source0.tex}.
\begin{theorem}\label{thm:nu-finite-theta}
Let $n=2m+1$ be odd, and define
\[
N_m(q,a):=\sum_{j=0}^{m} q^{2j+1}\frac{(aq,q/a;q^2)_j}{(-q;q^2)_{j+1}}.
\]
Then
\begin{equation}\label{eq:nu-finite-theta}
(1+q^n)\sum_{j=0}^{m} q^{2j+1}
\frac{(q^{1-n},q^{n+1};q^2)_j}{(-q;q^2)_{j+1}}
=
(-1)^m q^{(m+1)(2m+1)}
\sum_{r=-m}^{m} (-1)^r q^{-2r^2-r}.
\end{equation}
\end{theorem}

\begin{proof}[Proof of Theorem~\ref{thm:nu-finite-theta}]
Write
\begin{align*}
L_m(q)&:=(1+q^{2m+1})\sum_{j=0}^{m} q^{2j+1}
\frac{(q^{-2m},q^{2m+2};q^2)_j}{(-q;q^2)_{j+1}},\\
R_m(q)&:=(-1)^m q^{(m+1)(2m+1)}
\sum_{r=-m}^{m} (-1)^r q^{-2r^2-r}.
\end{align*}
Then~\eqref{eq:nu-finite-theta} is exactly the identity $L_m(q)=R_m(q)$.

For $m\ge 1$ and $0\le j\le m$, define
\begin{align*}
F_{m,j}(q)&:=(1+q^{2m+1})q^{2j+1}
\frac{(q^{-2m},q^{2m+2};q^2)_j}{(-q;q^2)_{j+1}},\\
G_{m,j}(q)&:=q^{2m+1}(1+q^{2m})
\frac{(q^{-2m},q^{2m};q^2)_j}{(-q;q^2)_j}.
\end{align*}
We claim that
\begin{equation}\label{eq:nu-telescope}
F_{m,j}(q)+q^{4m+1}F_{m-1,j}(q)=G_{m,j}(q)-G_{m,j+1}(q).
\end{equation}
Indeed,
\begin{align*}
F_{m,j}(q)+q^{4m+1}F_{m-1,j}(q)
&=
q^{2j+1}\frac{(q^{-2m},q^{2m};q^2)_j}{(-q;q^2)_{j+1}}\\
&\quad\times\Biggl(
\frac{(1+q^{2m+1})(1-q^{2m+2j})}{1-q^{2m}}\\
&\qquad\qquad
+\frac{q^{4m+1}(1+q^{2m-1})(1-q^{-2m+2j})}{1-q^{-2m}}
\Biggr).
\end{align*}
Since
\[
\frac{1}{1-q^{-2m}}=-\frac{q^{2m}}{1-q^{2m}},
\]
the factor in parentheses equals
\[
\begin{aligned}
\frac{(1+q^{2m+1})(1-q^{2m+2j})}{1-q^{2m}}
&- \frac{q^{6m+1}(1+q^{2m-1})(1-q^{-2m+2j})}{1-q^{2m}}\\
&= (1+q^{2m})\bigl(1+q^{2m+1}+q^{4m}-q^{2m+2j}\bigr).
\end{aligned}
\]
and therefore
\begin{equation}\label{eq:nu-left-common}
\begin{aligned}
F_{m,j}(q)+q^{4m+1}F_{m-1,j}(q)
&=
q^{2j+1}\frac{(q^{-2m},q^{2m};q^2)_j}{(-q;q^2)_{j+1}}\\
&\quad\times(1+q^{2m})\bigl(1+q^{2m+1}+q^{4m}-q^{2m+2j}\bigr).
\end{aligned}
\end{equation}

On the other hand,
\begin{align*}
G_{m,j}(q)-G_{m,j+1}(q)
&=
q^{2m+1}(1+q^{2m})
\frac{(q^{-2m},q^{2m};q^2)_j}{(-q;q^2)_j}\\
&\quad\times
\left(
1-\frac{(1-q^{-2m+2j})(1-q^{2m+2j})}{1+q^{2j+1}}
\right).
\end{align*}
Using $(-q;q^2)_{j+1}=(-q;q^2)_j(1+q^{2j+1})$, this becomes
\[
q^{2m+1}(1+q^{2m})
\frac{(q^{-2m},q^{2m};q^2)_j}{(-q;q^2)_{j+1}}
\Bigl((1+q^{2j+1})-(1-q^{-2m+2j})(1-q^{2m+2j})\Bigr).
\]
Now
\[
(1+q^{2j+1})-(1-q^{-2m+2j})(1-q^{2m+2j})
=
q^{-2m+2j}(1+q^{2m+1}+q^{4m}-q^{2m+2j}),
\]
so
\begin{equation}\label{eq:nu-right-common}
G_{m,j}(q)-G_{m,j+1}(q)
=
q^{2j+1}\frac{(q^{-2m},q^{2m};q^2)_j}{(-q;q^2)_{j+1}}
(1+q^{2m})(1+q^{2m+1}+q^{4m}-q^{2m+2j}).
\end{equation}
Comparing~\eqref{eq:nu-left-common} and~\eqref{eq:nu-right-common} proves~\eqref{eq:nu-telescope}.

Now $F_{m-1,m}(q)=0$ and $G_{m,m+1}(q)=0$. Summing~\eqref{eq:nu-telescope} over $j=0,\dots,m$, we obtain
\[
L_m(q)+q^{4m+1}L_{m-1}(q)
=
\sum_{j=0}^{m}\bigl(G_{m,j}(q)-G_{m,j+1}(q)\bigr)
=
G_{m,0}(q)
=
q^{2m+1}(1+q^{2m}).
\]
Hence
\begin{equation}\label{eq:nu-L-recurrence}
L_m(q)=-q^{4m+1}L_{m-1}(q)+q^{2m+1}(1+q^{2m}),
\qquad
L_0(q)=q.
\end{equation}

Now set
\[
\Psi_m(q):=\sum_{r=-m}^{m}(-1)^r q^{-2r^2-r}.
\]
Then
\[
\Psi_m(q)-\Psi_{m-1}(q)
=
(-1)^m\bigl(q^{-2m^2-m}+q^{-2m^2+m}\bigr)
=
(-1)^m q^{-2m^2-m}(1+q^{2m}).
\]
Therefore
\begin{align*}
R_m(q)+q^{4m+1}R_{m-1}(q)
&=
(-1)^m q^{(m+1)(2m+1)}\bigl(\Psi_m(q)-\Psi_{m-1}(q)\bigr)\\
&=
q^{2m+1}(1+q^{2m}).
\end{align*}
Together with $R_0(q)=q$, this shows that $R_m(q)$ satisfies the same recurrence~\eqref{eq:nu-L-recurrence} and the same initial value as $L_m(q)$. Thus $L_m(q)=R_m(q)$ for all $m\ge 0$.
\end{proof}

\end{document}
